% ─── CHAPTER 1 ────────────────────────────────────────────────────────────────
\chapter{Introduction}\label{INTRO}

\section{Motivation and Background}\label{INTRO_MOT}

Climate change is one of the defining challenges of the twenty-first century.
Understanding historical climate patterns --- temperature trends, precipitation
cycles, and the differences between regions --- is an essential first step for
researchers, educators, policymakers, and the general public alike. Yet access
to high-quality climate data remains, in practice, restricted to those with
specialist GIS or programming skills. Datasets such as
WorldClim~\cite{fick2017worldclim} provide kilometre-scale rasters of monthly
temperature and precipitation going back to 1960, but querying them
traditionally requires desktop software, custom scripts, or academic
subscriptions.

At the same time, the web platform has matured to the point where rich,
interactive data exploration can be delivered entirely in the browser, with no
installation and no registration. Modern JavaScript frameworks, vector-tile
mapping libraries, and reactive data-fetching primitives make it possible to
build tools that would previously have required a dedicated server-side
pipeline. The WorldClim group at the Universidad de Valladolid has developed a
Virtual Semantic API~\cite{scrapi} that exposes its raster archive as a
SPARQL-compatible REST service, opening the door to exactly this kind of
lightweight, client-only integration.

Despite these advances, no general-purpose, browser-based tool exists that
allows a non-specialist user to look up a city by name, inspect its monthly
climate profile, compare it to another city, explore how its climate has shifted
across historical periods, or visualise the spatial distribution of a climate
variable across an entire region --- all within a single, cohesive interface.
Climatica was designed to fill that gap.

\section{Objectives}\label{INTRO_OBJ}

The primary objective of this internship project was to design, implement, and
evaluate a web application that makes WorldClim historical climate data
accessible and explorable for non-specialist users, operating entirely in the
browser without requiring registration, local software installation, or GIS
expertise.

The specific objectives were as follows:

\begin{enumerate}
  \item \textbf{City climate profile.} Enable a user to search for any city in
    the world by name and view a complete monthly climate profile, including
    minimum temperature, maximum temperature, average temperature, and
    precipitation, together with a Walter-Lieth climatogram for rapid climate
    classification.

  \item \textbf{City comparison.} Allow side-by-side comparison of the climate
    profiles of two cities, with overlaid charts that make seasonal differences
    immediately visible.

  \item \textbf{Period comparison.} Expose the five historical time periods
    available in WorldClim (1960--1990, 1970--2000, 1980--2010, 1990--2020,
    and the long-term baseline) so that a user can observe trends in a city's
    climate across decades.

  \item \textbf{Regional heatmap.} Provide an interactive map on which the user
    can draw a bounding box or polygon and view a spatially resolved heatmap of
    any climate variable for any month, rendered directly from WorldClim raster
    data.

  \item \textbf{Front-end only architecture.} Implement the application as a
    single-page application that communicates directly with the WorldClim
    Virtual Semantic API and the Wikidata API, without developing a custom
    back-end. An early design considered wrapping the WorldClim API in a
    NestJS layer with a database and response caching, but this was set aside
    in favour of delivering a fully functional front-end within the available
    time.

  \item \textbf{Internationalisation.} Deliver the interface in at least two
    languages; the final implementation supports English, Spanish, and
    Ukrainian.
\end{enumerate}

\section{Methodology}\label{INTRO_MET}

The project followed an iterative, prototype-driven development process
organised into four broad phases.

\textbf{Phase 1 --- Research and scoping.} The first two weeks were devoted to
understanding the WorldClim Virtual Semantic API and the Wikidata knowledge
graph, surveying existing climate visualisation tools (see
Chapter~\ref{REL}), and agreeing the functional scope with the supervisor.
User stories were collected and translated into a set of functional and
non-functional requirements (Section~\ref{APP_REQ}).

\textbf{Phase 2 --- Architecture and technology selection.} Based on the
requirements, the technology stack was chosen (Section~\ref{APP_TECH}) and
the application architecture was defined (Section~\ref{APP_ARCH}). The key
architectural decision --- to build a front-end-only application that calls
the WorldClim and Wikidata APIs directly --- was validated against the
WorldClim API documentation and confirmed with the supervisor.

\textbf{Phase 3 --- Iterative implementation.} The four application pages were
implemented one at a time, starting with the City Climate page as it exercised
the core API integration and chart rendering. Each page was demonstrated to the
supervisor at fortnightly meetings, feedback was collected, and the next
iteration began. This phase constituted the majority of the 400 working hours.

\textbf{Phase 4 --- Evaluation and documentation.} A usability evaluation was
conducted with representative users (Section~\ref{EVAL}), the results were
analysed, and this report was written.

Throughout all phases the codebase was managed with Git and hosted on GitHub,
with a linear commit history that maps directly to the phases above.

\section{Structure of the Document}\label{INTRO_EST}

The remainder of this report is organised as follows.

Chapter~\ref{REL} reviews related work, covering existing climate data portals
and visualisation tools as well as the two external services on which Climatica
depends most heavily: the WorldClim dataset and its Virtual Semantic API, and
the Wikidata knowledge graph.

Chapter~\ref{APP} describes the application itself. It begins with the
functional and non-functional requirements, continues with the data model and
architectural decisions, details the technology choices and their rationale, and
concludes with a walkthrough of the four implemented pages.

Chapter~\ref{EVAL} presents the evaluation. It covers the demo scenarios used
to exercise the application, a usability study conducted with target users, and
a discussion of limitations and potential improvements.

Chapter~\ref{CONC} summarises the conclusions, reflects on whether the project
objectives were met, and outlines directions for future work.